\chapter{Introduction}

In this part we'll investigate 4 min-max theorems. In these theorems, a max property of a graph is often proved to be equal to a min property of the graph.

We should be aware that such min-max theorems are not seen everywhere. For example:

\begin{claim}
	\[\text{Max Matching} \leq \text{Min Vertex Cover}\]
\end{claim}
\begin{proof}
	Let $M$ be a matching and $C$ be a vertex cover. Then each edge in $M$ must have at least one endpoint in $C$. Since the edges in $M$ are disjoint, we have that $|M| \leq |C|$.
\end{proof}

The easiest counterexample to show that the above inequality can be strict is a triangle graph. The maximum matching has size 1, while the minimum vertex cover has size 2.

Here we first list the 4 min-max theorems and prove the $\leq$ direction. Then we will prove the tightness in the next few chapters.

\section{Menger's Theorem}

\begin{claim}[Menger's Theorem]
	In a unweighted graph $G$, the maximum number of edge-disjoint $s$-$t$ paths $\le$ the minimum $s$ - $t$ cut (i.e. the minimum number of edges whose removal disconnects $s$ and $t$).
\end{claim}
\begin{proof}
	Let $P_1, P_2, \ldots, P_k$ be edge-disjoint $s$-$t$ paths. Let $C$ be an $s$-$t$ cut. Then each path $P_i$ must contain at least one edge in $C$. Since the paths are edge-disjoint, we have that $k \leq |C|$.
\end{proof}

\section{Nash-Williams Theorem}

\begin{claim}[Nash-William Theorem]
	The maximum number of edge-disjoint spanning trees $\le \min \lfloor \frac{|E(P)|}{|P|-1} \rfloor$ over all partition $P$ of the graph.
\end{claim}
\begin{proof}
	Let $T_1, T_2, \ldots, T_k$ be edge-disjoint spanning trees. Let $P$ be a partition of the graph. Then each spanning tree $T_i$ must contain at least $|P|-1$ edges in $E(P)$ as a spanning tree must touch all the components of the partition $P$. Since the spanning trees are edge-disjoint, we have that $k(|P|-1) \leq |E(P)|$, which implies that $k \leq \lfloor \frac{|E(P)|}{|P|-1} \rfloor$.
\end{proof}

\begin{note}
	Intuitively we want to find an equality between the maximum number of disjoint spanning trees and the minimum global cut. However one can find the counterexample of it in $C_4$, with the maximum number of disjoint spanning trees being 1, while the minimum global cut is 2.
\end{note}

\section{Edmonds' Theorem}

\begin{definition}
	A $s$\textbf{-arborescence} in a directed graph is a spanning tree rooted at $s$ such that all vertices are reachable from $s$.
\end{definition}
\begin{claim}[Rooted Connectivity Theorem]
	The maximum number of edge-disjoint $s$-arborescences $\le$ the minimum $s$ cut (i.e. the minimum number of edges whose removal disconnects $s$ and any vertex).
\end{claim}

We won't be focusing on this theorem. Instead, we want to prove a weaker version in the next few chapters:

\begin{definition}
	A $k$ \textbf{bounded in-degree spanning tree} is a spanning tree packing in a directed graph with:
	\begin{itemize}
		\item $k$ edge-disjoint spanning trees $T_1, T_2, \ldots, T_k$.
		\item Let $T = \bigcup_{i=1}^k T_i$. Then for all $t \neq s, \indeg_T(t) = k$.
	\end{itemize}
\end{definition}

One can verify that the maximum number of edge-disjoint $s$-arborescences is at most the maximum number of $k$ bounded in-degree spanning trees.

\begin{claim}
	The maximum number of $k$ bounded in-degree spanning trees $\le$ the minimum $s$ cut.
\end{claim}
\begin{proof}
	Let the minimum $s$ cut separate a set $X$ from $V\setminus X$, and $s\in V\setminus X$. As every spanning tree can contain at most $|X|-1$ edges in the subgraph induced by $X$, we have
	\[|E_T(X)|\le k(|X|-1)\]
	Also, as every vertex in $X$ have in degree $k$, we have
	\[ \sum_{v\in X} \indeg_T(v) = k|X| \]
	Combining the two inequalities, we have
	\[ |E(V\setminus X, X)| = k|X| - |E_T(X)| \ge k|X| - k(|X|-1) = k \]
	which implies that the minimum $s$ cut is at least $k$.
\end{proof}

\section{Max-Dicut Packing Min-Dijoin Theorem}

\begin{definition}
	A \textbf{dicut} is the edge set $\delta^{\text{out}}(X)$ for any set $X$ that $\delta^{\text{in}}(X) = \emptyset$.

	A \textbf{dijoin} is a set of edges that, after contraction, makes the graph strongly connected (i.e. every vertex can reach every other vertex).
\end{definition}

\begin{claim}[Max-Dicut Packing Min-Dijoin Theorem]
	The maximum number of edge-disjoint dicuts $\le$ the minimum dijoin.
\end{claim}
\begin{proof}
	Let the maximum number of edge-disjoint dicuts be $k$. Let $D_1, D_2, \ldots, D_k$ be the edge-disjoint dicuts. Let $J$ be a dijoin. Then each dicut $D_i$ must contain at least one edge in $J$. Since the dicuts are edge-disjoint, we have that $k \leq |J|$.
\end{proof}
